%
\section{\crowsFoot{E}{O1}{F}{OM}}%
%
\begin{frame}[t]%
\frametitle{\crowsFoot{E}{O1}{F}{OM}}%
\begin{itemize}%
\item Es gibt zwei Entitätstypen~E und~F.%
\item<2-> Jede Entität vom Typ~E kann mit keiner, einer, oder mehreren Entitäten vom Typ~F verbunden sein.
\item<3-> Jede Entität vom Typ~F ist mit keiner oder einer Entität vom Typ~E verbunden.
\item<4-> Beispiel:\uncover<5->{%
\begin{itemize}%
%
\item Eine Nachricht kann entweder eine neue Nachricht sein, oder die Antwort auf eine existierende Nachricht.\uncover<6->{ %
Es kann keine, eine, oder beliebig viele Antworten auf jede Nachricht geben.}%
\end{itemize}%
}%
\end{itemize}%
%
\locateGraphic{4-}{width=0.8\paperwidth}{graphics/EF}{0.1}{0.6}%
\end{frame}%
%
\begin{frame}%
\frametitle{Lösungen}%
\begin{itemize}%
\item Wir brauchen eine Tabelle~\sqlil{e} für die Entitäten vom Typ~E.%
%
\item<2->  Wir brauchen auch eine Tabelle~\sqlil{f} für die Entitäten vom Typ~F.%
%
\item<3-> Beide haben ihre Ersatz-Primärschlüssel~\sqlil{eid} und~\sqlil{fid}.%
%
\item<4-> Wir fügen auch wieder die Beispielattribute~\sqlil{x} und~\sqlil{y} hinzu.%
%
\item<5-> Wieder gibt es zwei mögliche Lösungen für dieses Szenario\only<-5>{.}\uncover<6->{:%
\begin{enumerate}%
\item Wir können drei Tabellen nehmen, wobei die dritte Tabelle wieder die \inQuotes{Verbindungen} speichert\only<-6>{.}\uncover<7->{ oder}%
\item<7-> wir können zwei Tabellen nehmen, wobei die \inQuotes{Verbindungsdaten} diesmal in Tabelle~\sqlil{f} gehen müssen.%
\end{enumerate}%
}%
%
\item<8-> Welche Lösung besser ist, hängt wieder mit der Anzahl der verbundenen Datensätze zusammen.
\end{itemize}%
\end{frame}%
%

\begin{frame}[t]%
\frametitle{Mit Drei Tabellen}%
\parbox{0.435\linewidth}{\small%
\begin{itemize}%
%
\only<-4>{%
\item Wir benutzen eine dritte Tabelle~\sqlil{relate_e_and_f}, die die Entitäten vom Typ~E mit denen vom Typ~F in Verbindung setzt.%
}%
\only<-5>{%
\item<2-> In der Tabelle gibt es zwei Spalten.%
}%
\only<-6>{%
\item<3-> Die erste, \sqlil{fkeid}, speichert die Werte der Primärschlüssel~\sqlil{eid} der E-Entitäten.%
}%
\only<-7>{%
\item<4-> Die zweite, \sqlil{fkfid}, speichert die Werte der Primärschlüssel~\sqlil{fid} der F-Entitäten.%
}%
%
\only<-7>{%
\item<5-> Da wir nur Paare speichern, die existieren, sind beide Spalten~\sqlil{NOT NULL} und haben \sqlil{REFERENCES}-Einschränkungen auf ihre entsprechenden Fremdschlüssel.%
}%
%
\item<6-> Da jede Entität vom Typ~F nur mit höchstens einer Entität vom Type~E verbunden sein kann, gibt es eine \sqlil{UNIQUE}-Einschränkung auf Spalte~\sqlil{fkfid}.%
%
\item<7-> Wir machen sie auch zum Primärschlüssel der Tabelle~\sqlil{relate_e_and_f}.%
%
\item<8-> Sie muss der Primärschlüssel sein, denn die Werte in Spalte~\sqlil{fkeid} müssen nicht einmalig sein {\dots} eine Entität vom Typ~E kann mit vielen Entitäten vom Typ~F verbunden sein.%
%
\end{itemize}%
}%
\gitLoadAndExecSQL{EF_1_tables}{}{conceptualToRelational}{EF_1_tables.sql}{relationships}{}{}%
\listingSQLandOutput{}{EF_1_tables}{}{0.47}{0.1}{0.529}{0.87}
\end{frame}%
%
\begin{frame}[t]%
\frametitle{Befüllen mit Daten}%
\parbox{0.435\linewidth}{\small%
\begin{itemize}%
%
\item Wir füllen nun einige Daten in die Tabellen.%
%
\item<2-> Weil beide Enden der Beziehung optional sind, können wir erstmal die Tabellen~\sqlil{e} und~\sqlil{f} befüllen.%
%
\item<3-> Dann erstellen wir die Beziehungen, in dem wir Zeilen in die Tabelle~\sqlil{relate_e_and_f} einfügen.
\end{itemize}%
}%
\gitLoadAndExecSQL{EF_1_insert_and_select}{}{conceptualToRelational}{EF_1_insert_and_select.sql}{relationships}{}{}%
\listingSQLandOutput{}{EF_1_insert_and_select}{}{0.47}{0.1}{0.529}{0.87}
\end{frame}%
%
\begin{frame}[t]%
\frametitle{Der Inhalt der Drei Tabellen}%
%
\gitExec{cdtrmTableEO}{\databasesCodeRepo}{conceptualToRelational}{../_scripts_/db_table_to_latex_table.sh relationships e eid;x}%
\gitExec{cdtrmTableFO}{\databasesCodeRepo}{conceptualToRelational}{../_scripts_/db_table_to_latex_table.sh relationships f fid;y}%
\gitExec{cdtrmTableREF}{\databasesCodeRepo}{conceptualToRelational}{../_scripts_/db_table_to_latex_table.sh relationships relate_e_and_f fkeid;fkfid}%
%
\vfill%
%
\begin{itemize}%
\item Dies sind die Daten in den drei Tabellen.%
\end{itemize}%
\vfill%
%
\begin{center}%
\strut\hfill\strut%
\centering\input{\gitFile{cdtrmTableEO}}%
\strut\hfill\strut%
\centering\input{\gitFile{cdtrmTableFO}}%
\strut\hfill\strut%
\centering\input{\gitFile{cdtrmTableREF}}%
\strut\hfill\strut%
\end{center}%
%
\hfill%
\end{frame}%
%%
\begin{frame}%
\frametitle{Probleme mit der Drei-Tabellen-Methode}%
\begin{itemize}%
\item Dieser Ansatz wird \DEzB\ in \bracketCite{S2024D:MEDTRDM} empfohlen.%
%
\item<2-> Er hat aber wieder genau die selben Probleme, wie der drei-Tabellen-Ansatz für die Entitätstypen~A und~B.
%
\item<3-> Der Ansatz hat nur dann Sinn, wenn es relativ wenig Paare von verbundenen E~und F~Entitäten gibt.%
%
\item<4-> Wenn fast alle Entitäten vom Typ~F mit einer Entität vom Typ~E verbunden sind, dann ist die Tabelle~\sqlil{relate_e_and_f} fast genauso groß wie Tabelle~\sqlil{f}.%
%
\item<5-> Anstatt einfach die Schlüssel der verbundenen E-Entitäten mit in Tabelle~\sqlil{f} zu speichern, haben wir dann eine weitere Tabelle, die diese Schlüssel \emph{und} die Primärschlüssel von Tabelle~\sqlil{f} speichert.%
%
\item<6-> Und wir brauchen auch zwei \sqlil{INNER JOIN}s.
\end{itemize}%
\end{frame}%
%
\begin{frame}%
\frametitle{Zwei-Tabellen Lösung}%
\parbox{0.435\linewidth}{\small%
\begin{itemize}%
%
\only<-5>{%
\item Probieren wir also eine Lösung mit zwei Tabellen.%
}%
%
\only<-6>{%
\item<2-> Dazu löschen wir erstmal wieder unsere Tabellen.%
}%
%
\only<-7>{%
\item<3-> Diesmal müssen wir eine Fremdschlüssel-Spalte~\sqlil{fkeid} zur Tabelle~\sqlil{f} hinzufügen.%
}%
%
\item<4-> Diese Spalte bekommt eine \sqlil{REFERENCES}-Einschränkung, die auf Spalte~\sqlil{eid} der Tabelle~\sqlil{e} zeigt.%
%
\item<5-> Ein Wert in dieser Spalte darf \sqlil{NULL} sein, denn nicht alle Zeilen in Tabelle~\sqlil{f} müssen mit Zeilen in Tabelle~\sqlil{e} verbunden sein.%
%
\item<6-> Anders als in der \crowsFoot{A}{O1}{B}{O1}-Situation darf hier keine \sqlil{UNIQUE}-Einschränkung verwendet werden.%
%
\item<7-> Jede Zeile von Tabelle~\sqlil{e} kann mit mehreren Zeilen in Tabelle~\sqlil{f} verbunden sein.%
%
\item<8-> Die Primärschlüsselwerte von Tabelle~\sqlil{e} dürfen mehrmals als Fremdschlüsselwerte in Spalte~\sqlil{fkeid} in Tabelle~\sqlil{f} auftauchen.%
%
\end{itemize}%
}%
\gitLoadAndExecSQL{EF_cleanup}{}{conceptualToRelational}{EF_cleanup.sql}{relationships}{}{}%
\listingSQLandOutput{2}{EF_cleanup}{}{0.47}{0.1}{0.529}{0.87}
\gitLoadAndExecSQL{EF_2_tables}{}{conceptualToRelational}{EF_2_tables.sql}{relationships}{}{}%
\listingSQLandOutput{3-}{EF_2_tables}{}{0.47}{0.1}{0.529}{0.87}
\end{frame}%
%
\begin{frame}[t]%
\frametitle{Befüllen mit Daten}%
\parbox{0.435\linewidth}{\small%
\begin{itemize}%
\only<-5>{%
\item Wir speichern nun Daten in die beiden Tabellen.%
}%
%
\only<-6>{%
\item<2-> Wir fügen zuerst Zeilen in die Tabelle~\sqlil{e} ein.%
}%
%
\only<-7>{%
\item<3-> Dann fügen wir Zeilen in die Tabelle~\sqlil{f} ein.%
}%
%
\item<4-> Für jede dieser Zeilen können wir einen gültigen Fremdschlüsselwert~\sqlil{fkeid} auf eine Zeile in Tabelle~\sqlil{e}~(identifiziert vom Primärschlüssel~\sqlil{eid}) angeben.%
%
\item<5-> Wenn wir das machen, dann ist die neue Zeile in Tabelle~\sqlil{f} mit einer existierenden Zeile in Tabelle~\sqlil{e} verbunden.%
%
\item<6-> Wir können auch \sqlil{NULL} als \sqlil{fkeid} speichern.%
%
\item<7-> Dann ist die Zeile in Tabelle~\sqlil{f} nicht mit einer Zeile in Tabelle~\sqlil{e} verbunden.%
%
\item<8-> Mit einem einzelnen \sqlil{INNER JOIN} können wir die Daten wieder zusammensetzen.%
\end{itemize}%
}%
\gitLoadAndExecSQL{EF_2_insert_and_select}{}{conceptualToRelational}{EF_2_insert_and_select.sql}{relationships}{}{}%
\listingSQLandOutput{}{EF_2_insert_and_select}{}{0.47}{0.1}{0.529}{0.87}
\end{frame}%
%
\begin{frame}[t]%
\frametitle{Der Inhalt der Zwei Tabellen}%
%
\gitExec{cdtrmTableET}{\databasesCodeRepo}{conceptualToRelational}{../_scripts_/db_table_to_latex_table.sh relationships e eid;x}%
\gitExec{cdtrmTableFT}{\databasesCodeRepo}{conceptualToRelational}{../_scripts_/db_table_to_latex_table.sh relationships f fid;fkeid;y}%
%
\vfill%
%
\begin{itemize}%
\item Dies sind die Daten in den zwei Tabellen.%
\end{itemize}%
\vfill%
%
\begin{center}%
\strut\hfill\strut%
\centering\input{\gitFile{cdtrmTableET}}%
\strut\hfill\strut%
\centering\input{\gitFile{cdtrmTableFT}}%
\strut\hfill\strut%
\end{center}%
%
\hfill%
\end{frame}%
%
